% Options for packages loaded elsewhere
\PassOptionsToPackage{unicode}{hyperref}
\PassOptionsToPackage{hyphens}{url}
%
\documentclass[
]{article}
\usepackage{amsmath,amssymb}
\usepackage{iftex}
\ifPDFTeX
  \usepackage[T1]{fontenc}
  \usepackage[utf8]{inputenc}
  \usepackage{textcomp} % provide euro and other symbols
\else % if luatex or xetex
  \usepackage{unicode-math} % this also loads fontspec
  \defaultfontfeatures{Scale=MatchLowercase}
  \defaultfontfeatures[\rmfamily]{Ligatures=TeX,Scale=1}
\fi
\usepackage{lmodern}
\ifPDFTeX\else
  % xetex/luatex font selection
\fi
% Use upquote if available, for straight quotes in verbatim environments
\IfFileExists{upquote.sty}{\usepackage{upquote}}{}
\IfFileExists{microtype.sty}{% use microtype if available
  \usepackage[]{microtype}
  \UseMicrotypeSet[protrusion]{basicmath} % disable protrusion for tt fonts
}{}
\makeatletter
\@ifundefined{KOMAClassName}{% if non-KOMA class
  \IfFileExists{parskip.sty}{%
    \usepackage{parskip}
  }{% else
    \setlength{\parindent}{0pt}
    \setlength{\parskip}{6pt plus 2pt minus 1pt}}
}{% if KOMA class
  \KOMAoptions{parskip=half}}
\makeatother
\usepackage{xcolor}
\setlength{\emergencystretch}{3em} % prevent overfull lines
\providecommand{\tightlist}{%
  \setlength{\itemsep}{0pt}\setlength{\parskip}{0pt}}
\setcounter{secnumdepth}{-\maxdimen} % remove section numbering
\ifLuaTeX
  \usepackage{selnolig}  % disable illegal ligatures
\fi
\IfFileExists{bookmark.sty}{\usepackage{bookmark}}{\usepackage{hyperref}}
\IfFileExists{xurl.sty}{\usepackage{xurl}}{} % add URL line breaks if available
\urlstyle{same}
\hypersetup{
  hidelinks,
  pdfcreator={LaTeX via pandoc}}

\author{}
\date{}

\begin{document}

\hypertarget{guida}{%
\section{Guida}\label{guida}}

\hypertarget{sintesi-finale-per-il-tuo-studio-attivo}{%
\subsubsection{\texorpdfstring{\textbf{SINTESI FINALE PER IL TUO STUDIO ATTIVO}}{SINTESI FINALE PER IL TUO STUDIO ATTIVO}}\label{sintesi-finale-per-il-tuo-studio-attivo}}

Hai ora tra le mani una spiegazione integrata e discorsiva. Per trasformare questo in studio attivo, ti consiglio di fare così quando lo riscriverai:

\begin{enumerate}
\def\labelenumi{\arabic{enumi}.}
\tightlist
\item
  \textbf{Ripercorri il flusso logico:} Inizia sempre dalla griglia di numeri (0-255) e dalla formula \texttt{g\ =\ T(f)}.
\item
  \textbf{Disegna i grafici LUT:} Per ogni operazione (negativo, logaritmo, gamma\textgreater1, gamma\textless1, binarizzazione), prendi un foglio e \textbf{disegna a mano} l'asse X e l'asse Y. Poi traccia la curva caratteristica. Questo è l'esercizio più importante.
\item
  \textbf{Fai esempi numerici:} Prendi tre valori: un pixel scuro (30), uno medio (128), uno chiaro (220). Applica mentalmente le regole e chiediti: ``Questo operatore lo schiarisce o lo scurisce? Di molto o di poco?''.
\item
  \textbf{Collega la formula al grafico:} Quando vedi \texttt{s\ =\ c\ *\ r\^{}γ}, non memorizzarla. Pensa: ``Ah, γ è l'esponente. Se è piccolo (\textless1), la curva sul grafico sarà convessa verso l'alto, quindi schiarirà le ombre''.
\item
  \textbf{Rispondi alle domande nel testo:} Dove ti ho lasciato una domanda (``Domanda per te\ldots{}''), fermati e prova a ragionare con la tua testa prima di cercare la risposta.
\end{enumerate}

Il tuo obiettivo non è ricordare a memoria le formule, ma essere in grado di \textbf{guardare una curva LUT e prevedere l'effetto sull'immagine}, e viceversa, di \textbf{descrivere con parole tue l'effetto che vuoi ottenere e sapere che tipo di curva ti serve}.

Buon lavoro di riscrittura! È nel momento in cui rielaborerai questi concetti con le tue parole e i tuoi disegni che diventeranno davvero tuoi.

\end{document}
